\noindent\textit{Sửa từ đề cũ số 15. Biên soạn: Trang.}

\begin{bt}[2,0 điểm]
Tính hợp lí (nếu có thể):
\begin{listEX}[1]
\item $\dfrac{1}{5}\cdot\left(\dfrac{4}{3}\right)^{2}-\dfrac{1}{5}:\dfrac{9}{11}$
\item $\left(-\dfrac{3}{4}+\dfrac{2}{7}\right):\dfrac{2}{3}-\left(\dfrac{1}{4}-\dfrac{5}{7}\right)\cdot\dfrac{3}{2}$
\item $\sqrt{0{,}49}\cdot\left(\dfrac{1}{\sqrt{81}}+\dfrac{1}{3}-1\right)+(0{,}4)^{5}\cdot\left(\dfrac{5}{2}\right)^{5}$
\item $\dfrac{2^{10}\cdot 3^{8}}{6^{6}\cdot 4^{2}}$
\end{listEX}
\loigiai{
a) $\dfrac{1}{9}$.\quad
b) $0$.\quad
c) $\dfrac{11}{18}$.\quad
d) $9$.
}
\end{bt}

\begin{bt}[2,0 điểm]
Tìm $x$, biết:
\begin{listEX}[2]
\item $-\dfrac{5}{8}-\dfrac{41}{8}x=\dfrac{2}{5}$
\item $\left(-\dfrac{3}{4}\right)^{3x-1}=-\dfrac{27}{64}$
\item $\left|x+\dfrac{1}{5}\right|+\left|\dfrac{7}{2}+2x\right|=5x$
\item $2^{x+2}-2^{x}=48$
\end{listEX}
\loigiai{
a) $x=-\dfrac{1}{5}$.\\
b) $x=\dfrac{4}{3}$.\\
c) $x=\dfrac{37}{20}$. Vì vế phải không âm nên $x\ge 0$; khi đó cả hai biểu thức trong dấu giá trị tuyệt đối đều dương.\\
d) $x=4$.
}
\end{bt}

\begin{bt}[1,0 điểm]
Một người đi từ $A$ đến $B$ với vận tốc $12\,\mathrm{km/h}$, hết $2$ giờ $30$ phút. Lúc về, người đó đi đúng quãng đường $AB$ với vận tốc $15\,\mathrm{km/h}$.
\begin{listEX}[1]
\item Tính thời gian đi từ $B$ về $A$.
\item Tính vận tốc trung bình của cả hành trình đi và về.
\end{listEX}
\loigiai{
a) $2$ giờ.\\
b) $\dfrac{40}{3}\,\mathrm{km/h}$.\\
Quãng đường $AB$ dài $30\,\mathrm{km}$. Tổng thời gian là $4{,}5$ giờ, tổng quãng đường là $60\,\mathrm{km}$. Vận tốc trung bình không bằng trung bình cộng của $12\,\mathrm{km/h}$ và $15\,\mathrm{km/h}$.
}
\end{bt}

\begin{bt}[1,0 điểm]
Tính tổng số học sinh lớp $7$A và lớp $7$B của một trường, biết số học sinh lớp $7$A bằng $\dfrac{4}{5}$ số học sinh lớp $7$B. Nếu chuyển $8$ học sinh từ lớp $7$A sang lớp $7$B thì số học sinh lớp $7$A bằng $\dfrac{1}{2}$ số học sinh lớp $7$B.
\loigiai{
$72$ học sinh.\\
Gọi số học sinh lớp $7$B lúc đầu là $5k$ ($k$ nguyên dương), số học sinh lớp $7$A là $4k$. Ta có $4k-8=\dfrac{1}{2}(5k+8)$. Suy ra $k=8$, nên $7$A có $32$ em và $7$B có $40$ em.
}
\end{bt}

\begin{bt}[1,5 điểm]
Cho hình vẽ sau, biết $HD\parallel GE$ và $GE\parallel IF$.
\begin{listEX}[1]
\item Hai đoạn thẳng $HD$ và $IF$ có song song với nhau không? Vì sao?
\item Cho $\widehat{HAB}=65^{\circ}$. Tính số đo $x$ của $\widehat{ABE}$ và số đo $y$ của $\widehat{BCI}$.
\end{listEX}
\begin{center}
\begin{tikzpicture}[scale=1.15,line cap=round]
	\coordinate (H) at (0,2.4);
	\coordinate (D) at (4.2,2.4);
	\coordinate (G) at (0,1.2);
	\coordinate (E) at (4.2,1.2);
	\coordinate (I) at (0,0);
	\coordinate (F) at (4.2,0);
	\coordinate (T) at (2.85,3.12);
	\coordinate (A) at (2.55,2.4);
	\coordinate (B) at (2.05,1.2);
	\coordinate (C) at (1.55,0);
	\coordinate (K) at (1.25,-0.72);
	\draw (H) -- (D);
	\draw (G) -- (E);
	\draw (I) -- (F);
	\draw (T) -- (K);
	\node[left] at (H) {$H$};
	\node[right] at (D) {$D$};
	\node[left] at (G) {$G$};
	\node[right] at (E) {$E$};
	\node[left] at (I) {$I$};
	\node[right] at (F) {$F$};
	\node[above left] at (A) {$A$};
	\node[below left] at (B) {$B$};
	\node[below] at (C) {$C$};
	\node[above] at (T) {$T$};
	\node[below] at (K) {$K$};
	\draw ($(A)+(180:0.55)$) arc (180:247:0.55);
	\node at ($(A)+(214:0.88)$) {$65^{\circ}$};
	\draw ($(B)+(0:0.48)$) arc (0:67:0.48);
	\node at ($(B)+(34:0.78)$) {$x$};
	\draw ($(C)+(67:0.5)$) arc (67:180:0.5);
	\node at ($(C)+(124:0.82)$) {$y$};
\end{tikzpicture}
\end{center}
\loigiai{
a) Có. Vì $HD\parallel GE$ và $GE\parallel IF$ nên $HD\parallel IF$.\\
b) $x=65^{\circ}$, $y=115^{\circ}$.\\
$\widehat{HAB}$ và $\widehat{ABG}$ là hai góc trong cùng phía, nên $\widehat{ABG}=115^{\circ}$. Góc này kề bù với $\widehat{ABE}$, vậy $x=65^{\circ}$. Góc tương ứng của $x$ tại $C$ là $\widehat{BCF}=65^{\circ}$, suy ra $y=180^{\circ}-65^{\circ}=115^{\circ}$.
}
\end{bt}

\begin{bt}[2,0 điểm]
Cho tam giác đều $ABC$. Dựng tam giác đều $BCD$ sao cho $D$ và $A$ nằm khác phía đối với đường thẳng $BC$. Gọi $H$ là trung điểm của $BC$.
\begin{listEX}[1]
\item Chứng minh $\triangle ABD$ cân và tính $\widehat{ABD}$, $\widehat{BAD}$.
\item Chứng minh $AH\perp BC$.
\item Chứng minh ba điểm $A$, $H$, $D$ thẳng hàng.
\end{listEX}
\begin{center}
\begin{tikzpicture}[scale=1.05,line join=round]
	\coordinate (B) at (-1.5,0);
	\coordinate (C) at (1.5,0);
	\coordinate (A) at (0,2.598);
	\coordinate (D) at (0,-2.598);
	\coordinate (H) at (0,0);
	\draw (A)--(B)--(C)--cycle;
	\draw (B)--(D)--(C);
	\draw (A)--(D);
	\node[above] at (A) {$A$};
	\node[left] at (B) {$B$};
	\node[right] at (C) {$C$};
	\node[below] at (D) {$D$};
	\node[above right] at (H) {$H$};
\end{tikzpicture}
\end{center}
\loigiai{
a) $AB=BC=BD$ nên $AB=BD$. $\widehat{ABD}=\widehat{ABC}+\widehat{CBD}=120^{\circ}$. Tam giác $ABD$ cân tại $B$, hai góc đáy bằng $30^{\circ}$. Vậy $\widehat{BAD}=30^{\circ}$.\\
b) Tam giác $ABC$ đều, $H$ là trung điểm của $BC$ nên $AH$ là đường trung tuyến, đồng thời là đường cao. Vậy $AH\perp BC$.\\
c) $A$, $H$, $D$ thẳng hàng.\\
\textit{Hướng dẫn.} Trong tam giác đều $ABC$, $\widehat{BAH}=30^{\circ}$. Góc này bằng $\widehat{BAD}$ nên tia $AH$ trùng tia $AD$.
}
\end{bt}

\begin{bt}[0,5 điểm]
Không dùng máy tính cầm tay, hãy so sánh $\sqrt{35}+\sqrt{99}$ và $16$.
\loigiai{
$\sqrt{35}+\sqrt{99}<16$.\\
\textit{Hướng dẫn.} $\sqrt{35}<\sqrt{36}=6$ và $\sqrt{99}<\sqrt{100}=10$.
}
\end{bt}
